\begin{thebibliography}{99}
\bibitem{fachantidis} Fachantidis, A., Taylor, M. E., \& Vlahavas, I. (2019). Learning to teach reinforcement learning agents. \emph{Machine Learning and Knowledge Extraction}, 1(1), 21--42. \url{https://doi.org/10.3390/make1010002}
\bibitem{najar} Najar, A., \& Chetouani, M. (2021). Reinforcement learning with human advice: A survey. \emph{Frontiers in Robotics and AI}, 8, 584075. \url{https://doi.org/10.3389/frobt.2021.584075}
\bibitem{chernova} Chernova, S., \& Veloso, M. (2009). Interactive policy learning through confidence-based autonomy. \emph{Journal of Artificial Intelligence Research}, 34, 1--25. \url{https://doi.org/10.1613/jair.2584}
\bibitem{zhu} Zhu, C., Cai, Y., Hu, S., Leung, H.-f., \& Chiu, D. K. W. (2023). Learning by reusing previous advice: A memory-based teacher--student framework. \emph{Autonomous Agents and Multi-Agent Systems}, 37, Article 14. \url{https://doi.org/10.1007/s10458-022-09595-1}
\bibitem{ilhan} \.{I}lhan, E., Gow, J., \& Perez-Liebana, D. (2021). Action advising with advice imitation in deep reinforcement learning. In \emph{Proceedings of the 20th International Conference on Autonomous Agents and Multiagent Systems} (pp. 629--637). \url{https://www.ifaamas.org/Proceedings/aamas2021/pdfs/p629.pdf}
\bibitem{act} Zhao, T. Z., Kumar, V., Levine, S., \& Finn, C. (2023). Learning fine-grained bimanual manipulation with low-cost hardware. In \emph{Proceedings of Robotics: Science and Systems XIX}. \url{https://doi.org/10.15607/RSS.2023.XIX.016}
\bibitem{rtc} Black, K., Galliker, M. Y., \& Levine, S. (2025). Real-time execution of action chunking flow policies. \emph{Advances in Neural Information Processing Systems}, 38. \url{https://doi.org/10.52202/085713-1122}
\bibitem{silva} da Silva, F. L., \& Costa, A. H. R. (2019). A survey on transfer learning for multiagent reinforcement learning systems. \emph{Journal of Artificial Intelligence Research}, 64, 645--703. \url{https://doi.org/10.1613/jair.1.11396}
\bibitem{agarwal} Agarwal, R., Schwarzer, M., Castro, P. S., Courville, A. C., \& Bellemare, M. G. (2021). Deep reinforcement learning at the edge of the statistical precipice. \emph{Advances in Neural Information Processing Systems}, 34, 29304--29320. 
\bibitem{pineau} Pineau, J., Vincent-Lamarre, P., Sinha, K., Larivi\`ere, V., Beygelzimer, A., d'Alch\'e-Buc, F., Fox, E., \& Larochelle, H. (2021). Improving reproducibility in machine learning research (a report from the NeurIPS 2019 reproducibility program). \emph{Journal of Machine Learning Research}, 22(164), 1--20. \url{https://jmlr.org/papers/v22/20-303.html}
\bibitem{metaworld} Yu, T., Quillen, D., He, Z., Julian, R., Hausman, K., Finn, C., \& Levine, S. (2020). Meta-World: A benchmark and evaluation for multi-task and meta reinforcement learning. \emph{Proceedings of Machine Learning Research}, 100, 1094--1100. \url{https://proceedings.mlr.press/v100/yu20a.html}

\bibitem{torrey} Torrey, L., \& Taylor, M. E. (2013). Teaching on a budget: Agents advising agents in reinforcement learning. In \emph{Proceedings of AAMAS} (pp. 1053--1060). \url{https://www.ifaamas.org/Proceedings/aamas2013/docs/p1053.pdf}
\bibitem{taylor} Taylor, M. E., \& Stone, P. (2009). Transfer learning for reinforcement learning domains: A survey. \emph{Journal of Machine Learning Research}, 10(56), 1633--1685. \url{https://jmlr.org/papers/v10/taylor09a.html}
\bibitem{griffith} Griffith, S., Subramanian, K., Scholz, J., Isbell, C. L., \& Thomaz, A. L. (2013). Policy shaping: Integrating human feedback with reinforcement learning. \emph{Advances in Neural Information Processing Systems}, 26. 
\bibitem{macglashan} MacGlashan, J., Ho, M. K., Loftin, R., Peng, B., Wang, G., Roberts, D. L., Taylor, M. E., \& Littman, M. L. (2017). Interactive learning from policy-dependent human feedback. \emph{Proceedings of Machine Learning Research}, 70, 2285--2294. \url{https://proceedings.mlr.press/v70/macglashan17a.html}
\bibitem{celemin} Celemin, C., \& Ruiz-del-Solar, J. (2019). An interactive framework for learning continuous actions policies based on corrective feedback. \emph{Journal of Intelligent \& Robotic Systems}, 95(1), 77--97. \url{https://doi.org/10.1007/s10846-018-0839-z}
\bibitem{dulac} Dulac-Arnold, G., Levine, N., Mankowitz, D. J., Li, J., Paduraru, C., Gowal, S., \& Hester, T. (2021). Challenges of real-world reinforcement learning: Definitions, benchmarks and analysis. \emph{Machine Learning}, 110(9), 2419--2468. \url{https://doi.org/10.1007/s10994-021-05961-4}
\bibitem{ramstedt} Ramstedt, S., \& Pal, C. (2019). Real-time reinforcement learning. \emph{Advances in Neural Information Processing Systems}, 32. 
\bibitem{randomdelay} Bouteiller, Y., Ramstedt, S., Beltrame, G., Pal, C., \& Binas, J. (2021). Reinforcement learning with random delays. In \emph{International Conference on Learning Representations}. \url{https://openreview.net/forum?id=QFYnKlBJYR}
\bibitem{diffusion} Chi, C., Feng, S., Du, Y., Xu, Z., Cousineau, E., Burchfiel, B. C., \& Song, S. (2023). Diffusion policy: Visuomotor policy learning via action diffusion. In \emph{Robotics: Science and Systems XIX}. \url{https://doi.org/10.15607/RSS.2023.XIX.026}
\bibitem{ross} Ross, S., Gordon, G., \& Bagnell, D. (2011). A reduction of imitation learning and structured prediction to no-regret online learning. \emph{Proceedings of Machine Learning Research}, 15, 627--635. \url{https://proceedings.mlr.press/v15/ross11a.html}
\bibitem{jiang} Jiang, N., \& Li, L. (2016). Doubly robust off-policy value evaluation for reinforcement learning. \emph{Proceedings of Machine Learning Research}, 48, 652--661. \url{https://proceedings.mlr.press/v48/jiang16.html}
\bibitem{thomas} Thomas, P., \& Brunskill, E. (2016). Data-efficient off-policy policy evaluation for reinforcement learning. \emph{Proceedings of Machine Learning Research}, 48, 2139--2148. \url{https://proceedings.mlr.press/v48/thomasa16.html}
\bibitem{swaminathan} Swaminathan, A., \& Joachims, T. (2015). Counterfactual risk minimization: Learning from logged bandit feedback. \emph{Proceedings of Machine Learning Research}, 37, 814--823. \url{https://proceedings.mlr.press/v37/swaminathan15.html}
\bibitem{dudik} Dud\'ik, M., Erhan, D., Langford, J., \& Li, L. (2014). Doubly robust policy evaluation and optimization. \emph{Statistical Science}, 29(4), 485--511. \url{https://doi.org/10.1214/14-STS500}
\bibitem{cawley} Cawley, G. C., \& Talbot, N. L. C. (2010). On over-fitting in model selection and subsequent selection bias in performance evaluation. \emph{Journal of Machine Learning Research}, 11(70), 2079--2107. \url{https://jmlr.org/papers/v11/cawley10a.html}
\bibitem{varma} Varma, S., \& Simon, R. (2006). Bias in error estimation when using cross-validation for model selection. \emph{BMC Bioinformatics}, 7, 91. \url{https://doi.org/10.1186/1471-2105-7-91}
\bibitem{henderson} Henderson, P., Islam, R., Bachman, P., Pineau, J., Precup, D., \& Meger, D. (2018). Deep reinforcement learning that matters. \emph{Proceedings of the AAAI Conference on Artificial Intelligence}, 32(1). \url{https://doi.org/10.1609/aaai.v32i1.11694}
\bibitem{fisher} Fisher, A., Rudin, C., \& Dominici, F. (2019). All models are wrong, but many are useful: Learning a variable's importance by studying an entire class of prediction models simultaneously. \emph{Journal of Machine Learning Research}, 20(177), 1--81. \url{https://jmlr.org/papers/v20/18-760.html}
\bibitem{hooker} Hooker, G., Mentch, L., \& Zhou, S. (2021). Unrestricted permutation forces extrapolation: Variable importance requires at least one more model, or there is no free variable importance. \emph{Statistics and Computing}, 31(6), 82. \url{https://doi.org/10.1007/s11222-021-10057-z}
\bibitem{strobl} Strobl, C., Boulesteix, A.-L., Kneib, T., Augustin, T., \& Zeileis, A. (2008). Conditional variable importance for random forests. \emph{BMC Bioinformatics}, 9, 307. \url{https://doi.org/10.1186/1471-2105-9-307}
\end{thebibliography}
